\documentclass[uplatex,dvipdfmx,11pt,a4paper]{bxjsarticle}
\usepackage{amsmath,amssymb}
\usepackage{url}
\usepackage[haranoaji]{pxchfon}

\begin{document}

\begin{center}
{\LARGE \textbf{Transmission Theory Lesson3 Report}}
\end{center}

\vspace{1em}

\begin{flushright}
2611193 Heizo Nakai
\end{flushright}

\section*{問1}

整合フィルタが受信フィルタとして最適である理由は、サンプリング時点でのSNRを最大化するように設計されているためである。受信信号には送信信号成分と雑音成分が含まれる。単に受信フィルタの帯域を狭めれば雑音成分を減らすことができるが、送信信号に必要な周波数成分まで失われてしまう。整合フィルタは、送信信号の波形に合わせて設計される。周波数領域では、送信信号の複素共役を用いることで、各周波数成分の位相を打ち消し、サンプリング時点で信号成分が同じ向きにそろって加算されるようにする。複素共役とは、幾何学的には複素平面上で実軸に対して反転する操作であり、信号の大きさは変えずに位相の向きを逆にする操作である。これにより、受信信号中に含まれる送信波形と同じ成分を強く取り出すことができる。一方、雑音は送信信号と無関係な成分であるため、整合フィルタを通しても、送信信号と同じ方向の成分は大きく取り出されるが、無関係な方向の雑音成分は相対的に抑えられる。したがって、整合フィルタはサンプリング時点でのSNRを最大化できるため、受信フィルタとして最適である。

\section*{問2}

2つのコンスタレーション点の距離が小さいほど、ビット誤り率は悪化する。通信路で加わる雑音により、受信された信号は本来のコンスタレーション点を中心として正規分布状に分布する。このとき、点間距離が小さいと、それぞれの確率密度関数（PDF）の曲線の重なりが大きくなる。重なった領域では、受信信号がどちらの信号点に対応するかの識別が難しくなり、判定境界を越えて隣の点として誤判定されやすい。したがって、点間距離が小さいほどBERは増加する。

\end{document}
